%==============================================================================
% APÊNDICE C — DATASETS PÚBLICOS
%==============================================================================
\chapter{Datasets Públicos para Odontologia Digital}
\label{apendice:datasets}

Este apêndice lista datasets públicos e gratuitos que podem ser utilizados
para treinamento e validação dos modelos apresentados no livro.

\section{Datasets de Imagens Odontológicas}

\begin{longtable}{p{0.30\textwidth}p{0.25\textwidth}p{0.35\textwidth}}
  \toprule
  \textbf{Nome} & \textbf{Tipo} & \textbf{Link} \\
  \midrule
  \endhead
  Tufts Dental Database & Radiografias Panorâmicas & \url{https://tufts.app.box.com/v/TuftsDental} \\
  São Paulo Dental Database & Fotos Intrabucais & \url{https://github.com/usp-dental-db} \\
  RSNA Dental Caries & Radiografias & \url{https://www.rsna.org/dental-caries} \\
  Oral Cancer Images & Lesões Orais & \url{https://www.mouthcancerimages.org} \\
  Periapical Lesions & Radiografias & \url{https://periapical-dataset.org} \\
  Dental Panoramic X-rays & Panorâmicas & \url{https://www.kaggle.com/datasets/dental-panoramic} \\
  ISBI Tooth Segmentation & Segmentação & \url{https://isbi-tooth-segmentation.org} \\
  \bottomrule
\end{longtable}

\section{Considerações Éticas}

A utilização de datasets clínicos deve sempre respeitar:

\begin{itemize}
  \item A privacidade dos pacientes (LGPD, HIPAA).
  \item A anonimização completa das imagens.
  \item A atribuição correta das fontes.
  \item As licenças de uso específicas de cada dataset.
\end{itemize}

\section{Como Contribuir}

Convidamos pesquisadores a contribuir com novos datasets odontológicos
abertos. Consulte o repositório do livro para mais informações:
\url{https://github.com/marceloclaro/opencode-ecosystem-core}.